% ---- pg-089 （书页 77） ----

奇异值分解的矩阵展开式用满维矩阵 $\bm{U}, \bm{V}, \bm{\Delta}$ 可写为

\[
\bm{A}_{(m \times k)} = \bm{U}_{(m \times m)}\bm{\Delta}_{(m \times k)}\bm{V}_{(k \times k)}'
\]

其中 $\bm{U}$ 的列为 $\bm{A}\bm{A}'$ 的 $m$ 个正交特征向量, $\bm{V}$ 的列为 $\bm{A}'\bm{A}$ 的 $k$ 个正交特征向量, $\bm{\Delta}$ 如结论 2A.15 中所示.

设

\[
\bm{A} = \begin{bmatrix} 3 & 1 & 1 \\ -1 & 3 & 1 \end{bmatrix}
\]

则

\[
\bm{A}\bm{A}' = \begin{bmatrix} 3 & 1 & 1 \\ -1 & 3 & 1 \end{bmatrix}\begin{bmatrix} 3 & -1 \\ 1 & 3 \\ 1 & 1 \end{bmatrix} = \begin{bmatrix} 11 & 1 \\ 1 & 11 \end{bmatrix}
\]

可以验证 $\bm{A}\bm{A}'$ 的特征值 $\gamma = \lambda^2$ 满足方程 $\gamma^2 - 22\gamma + 120 = (\gamma - 12)(\gamma - 10)$,因此,特征值是 $\gamma_1 = \lambda_1^2 = 12$ 和 $\gamma_2 = \lambda_2^2 = 10$,对应的特征向量分别是 $\bm{u}_1' = \left[\dfrac{1}{\sqrt{2}}, \dfrac{1}{\sqrt{2}}\right]$ 和 $\bm{u}_2' = \left[\dfrac{1}{\sqrt{2}}, -\dfrac{1}{\sqrt{2}}\right]$.

又

\[
\bm{A}'\bm{A} = \begin{bmatrix} 3 & -1 \\ 1 & 3 \\ 1 & 1 \end{bmatrix}\begin{bmatrix} 3 & 1 & 1 \\ -1 & 3 & 1 \end{bmatrix} = \begin{bmatrix} 10 & 0 & 2 \\ 0 & 10 & 4 \\ 2 & 4 & 2 \end{bmatrix}
\]

所以 $|\bm{A}'\bm{A} - \gamma\bm{I}| = -\gamma^3 - 22\gamma^2 - 120\gamma = -\gamma(\gamma - 12)(\gamma - 10)$,且特征值为 $\gamma_1 = \lambda_1^2 = 12$, $\gamma_2 = \lambda_2^2 = 10$ 及 $\gamma_3 = \lambda_3^2 = 0$.非零特征值与 $\bm{A}\bm{A}'$ 的非零特征值是相同的.用计算机算出特征向量为 $\bm{v}_1' = \left[\dfrac{1}{\sqrt{6}}, \dfrac{2}{\sqrt{6}}, \dfrac{1}{\sqrt{6}}\right]$, $\bm{v}_2' = \left[\dfrac{2}{\sqrt{5}}, -\dfrac{1}{\sqrt{5}}, 0\right]$ 及 $\bm{v}_3' = \left[\dfrac{1}{\sqrt{30}}, \dfrac{2}{\sqrt{30}}, -\dfrac{5}{\sqrt{30}}\right]$.对特征向量 $\bm{v}_1$ 和 $\bm{v}_2$ 可以进行检验：

\[
\bm{A}'\bm{A}\bm{v}_1 = \begin{bmatrix} 10 & 0 & 2 \\ 0 & 10 & 4 \\ 2 & 4 & 2 \end{bmatrix}\begin{bmatrix} 1 \\ 2 \\ 1 \end{bmatrix}\frac{1}{\sqrt{6}} \times \frac{1}{\sqrt{6}} = 12 \times \frac{1}{\sqrt{6}}\begin{bmatrix} 1 \\ 2 \\ 1 \end{bmatrix} = \lambda_1^2\bm{v}_1
\]
\[
\bm{A}'\bm{A}\bm{v}_2 = \begin{bmatrix} 10 & 0 & 2 \\ 0 & 10 & 4 \\ 2 & 4 & 2 \end{bmatrix}\begin{bmatrix} 2 \\ -1 \\ 0 \end{bmatrix}\frac{1}{\sqrt{5}} = 10 \times \frac{1}{\sqrt{5}}\begin{bmatrix} 2 \\ -1 \\ 0 \end{bmatrix} = \lambda_2^2\bm{v}_2
\]

选取 $\lambda_1 = \sqrt{12}$ 和 $\lambda_2 = \sqrt{10}$,我们求得 $\bm{A}$ 的奇异值分解是

\begin{equation*}
\begin{aligned}
\bm{A} &= \begin{bmatrix} 3 & 1 & 1 \\ -1 & 3 & 1 \end{bmatrix} \\
&= \sqrt{12}\begin{bmatrix} \dfrac{1}{\sqrt{2}} \\[2mm] \dfrac{1}{\sqrt{2}} \end{bmatrix}\left[\frac{1}{\sqrt{6}} \quad \frac{2}{\sqrt{6}} \quad \frac{1}{\sqrt{6}}\right] + \sqrt{10}\begin{bmatrix} \dfrac{1}{\sqrt{2}} \\[2mm] -\dfrac{1}{\sqrt{2}} \end{bmatrix}\left[\frac{2}{\sqrt{5}} \quad -\frac{1}{\sqrt{5}} \quad 0\right]
\end{aligned}
\end{equation*}

可以完成右边的计算来检验该等式.

奇异值分解与一个用低维矩阵逼近三角矩阵的结果有紧密的联系,该工作主要由 Eckart 和 Young(见文献[2])完成.如果做一个 $m \times k$ 矩阵 $\bm{A}$ 可以用具有相同维度但秩更低的矩阵 $\bm{B}$ 来逼近,差的平方和为
